\chapter{Introduction}
\label{ch:aims}

% Written 21 September 2026 from analysis/SPINE.md sections 1, 2 and 7b.

Language-model agents trained by reinforcement learning are increasingly
deployed where they will act alongside other agents that are themselves
learning.
Opponent shaping asks whether such an agent can be trained to influence how
a co-player learns, rather than only how to play against it.
Garc\'ia Segura et al.\ \cite{garciasegura2025} put that question to
language models: their shaper, trained on the return of a multi-episode
trial while a naive learner adapts against it, reaches the exploitative
outcome of the iterated Chicken game, earning 2.98 units per step against
its co-player's 1.01.

That result has two candidate causes.
The shaper is trained on a trial-level objective, which is the mechanism the
paper proposes.
It also learns an order of magnitude more slowly and updates five times less
often than its co-player, and in Chicken the player who commits wins, so an
agent that changes its policy slowly behaves more like a committed one.
The published design contains no learning-rate or update-schedule control:
its only shaping ablation enriches the naive learner's observations.
Which of the two causes produces the outcome is therefore open, and it
matters, because only one of them is a capability that transfers to settings
where commitment does not already pay.

\section{What answering the question requires}
\label{sec:aims-requires}

Two things, neither of which is present in the published comparison.

First, a control that differs from the shaper in the objective and in
nothing else.
A shaper does two things at once: it waits for a whole trial before
updating, and it lets credit cross episode boundaries within that trial.
Comparing it with an ordinary learner conflates the two.
The control used throughout this thesis is a \emph{trial-batched naive}
agent, which keeps the shaper's learning rate, clip, update schedule, batch
and prompt, and cuts credit at every episode boundary.
The difference between them is the objective and nothing else.

Second, an environment in which a co-player's policy is worth something to
the shaper whatever the shaper's own competence.
This is not automatic.
In the resource game this project began with, an agent that manages the
stock well keeps it alive against almost any partner, so a partner's
restraint is worth 2.0 units to a best responder and 0.0 from a full stock:
the apparent stake is a measure of the shaper's incompetence rather than of
what shaping buys.
The environment used here is built to remove that escape (Chapter~\ref{ch:env}).

\section{The environment and the design}
\label{sec:aims-design}

The experiment runs on a two-player stochastic common-pool resource with
integer stock, simultaneous takes of one or two units, stochastic regrowth
and a fixed 50-round horizon.
A single growth parameter switches it between two regimes.
At \(m=2\) no player can keep the pool alive alone and commitment loses to
reciprocity, so a shaper must induce restraint rather than out-commit its
co-player; a restrained co-player is worth 55.1 units.
At \(m=3\) a single player can sustain the pool and commitment wins, which
makes that regime a negative control: an exploitative outcome there is
explained by commitment, and a pre-registered rule says so before the data.
Both regimes are solved exactly, so every quantity an agent is compared
against is a number fixed before training.

The claim, the arms, the readouts and the decision rules were registered
before the runs, and the registration carries three gates that had to pass
first, two of which failed on earlier designs and are on record.
Forty-three training runs of 100 epochs were executed across nine arms at
\(m=2\) and three at \(m=3\).
What an agent acquired is read not from its co-play returns but from its
frozen policy, probed for 20 epochs against a scripted partner and assigned
a class.

\section{Findings}
\label{sec:aims-findings}

\begin{enumerate}
\item \textbf{The trial-level objective did not change what the co-learner
  acquired.} On the pre-registered arms the result is a null. It fails on the
  controls rather than on the shapers: two naive learners produced an
  unconditional restrainer on five of five seeds, and the shaping arms did
  not.
\item \textbf{Estimated by chaining advantages across the trial, the
  objective destroyed the resource.} Every such arm recorded survival 0.00
  on every seed, with both agents earning near the mutual-harvest value,
  while the controls sustained the pool and earned near mutual restraint.
\item \textbf{The estimator, not the objective, separates the arms.} An arm
  identical in every respect except how cross-episode credit is estimated
  sustained the resource and left an unconditional co-learner. It matched
  the controls; it did not beat them.
\item \textbf{The exploitative outcome appeared only where commitment
  explains it.} At \(m=3\) the ShapeLLM-style arm reproduced the published
  pattern, and so did an arm with no trial objective at all. The
  pre-registered negative control fires, and the outcome is not evidence of
  shaping.
\end{enumerate}

\section{Contributions}
\label{sec:aims-contributions}

\begin{itemize}
\item An exactly solved two-player stochastic commons in which the value of
  a co-player's restraint does not depend on the shaper's own competence,
  with a one-parameter negative control in which commitment pays.
\item An identification design for opponent shaping: the trial-batched
  control that separates the cross-episode objective from the update
  schedule and the timescale that accompany it.
\item A pre-registered null with its mechanism: in an environment verified
  learnable before training, the trial-level objective produced no change in
  the co-learner's acquired policy, and the measurements that say why.
\item The finding that credit assignment dominates the objective at this
  scale: two arms differing only in the estimator of the same objective
  differ in whether the resource survives at all.
\item A gate protocol that spends compute only after the environment is
  shown learnable, with three logged attempts, two of which stopped an
  experiment that would have produced an uninterpretable result.
\end{itemize}

\section{What this thesis does not claim}
\label{sec:aims-scope}

It does not show that opponent shaping fails for language-model agents.
One base model, one adapter rank, one environment family, 100 epochs and
three to five seeds bound every statement here.
It does not show that the trial-level objective is unlearnable, only that
this estimator of it did not learn in this setting.
It does not refute Garc\'ia Segura et al.: the environment class, the
hyperparameters and the readout all differ, and the reinterpretation offered
in Chapter~\ref{ch:discussion} concerns what that evidence form can
identify, not whether the reported result occurred.

\section{Organisation}
\label{sec:aims-organisation}

Chapter~\ref{ch:lit} places the work against opponent shaping, language-model
agents in social dilemmas, and the resource economics the environment is
drawn from.
Chapter~\ref{ch:env} defines the environment and solves it.
Chapter~\ref{ch:method} gives the agents, the three learning rules and the
readouts.
Chapter~\ref{ch:impl} covers the implementation and the tests.
Chapter~\ref{ch:design} is the pre-registration: arms, decision rules and
gates.
Chapter~\ref{ch:results} reports the runs, Chapter~\ref{ch:discussion}
interprets them, Chapter~\ref{ch:lim} states what bounds them, and
Chapter~\ref{ch:conc} concludes.
